\documentclass{beamer}

\usepackage[english]{babel}
\usepackage[utf8]{inputenc}
\usepackage[T1]{fontenc}
\usepackage{graphicx}
\usepackage{booktabs}
\usepackage{amsmath}
\usepackage{hyperref}

\mode<presentation>{
  \usetheme{default}
  \usecolortheme{default}
  \usefonttheme{default}
  \setbeamertemplate{caption}[numbered]
  \setbeamertemplate{navigation symbols}{}
}

% Notes remain in the source. To print them, uncomment the next line.
% \setbeameroption{show notes}
\newcommand{\code}[1]{\texttt{\detokenize{#1}}}

\title{DeepCASE Project}
\subtitle{Data input and graph construction}
\author{Allan Murillo}
\institute{TEC}
\date{October 11, 2026}

% Implementation sources:
% deepcase_gnn/preprocessing/preprocessor.py
% deepcase_gnn/preprocessing/lanl.py
% deepcase_gnn/context_builder/context_builder.py
% Dataset source: https://csr.lanl.gov/data/cyber1/
% This section describes the current DeepCASE-GNN implementation.
% All C10/C20/C30 records below are illustrative, not quoted dataset rows.

\begin{document}

\begin{frame}
  \titlepage
\end{frame}

\begin{frame}{Outline}
  \tableofcontents
\end{frame}

\section{Cross-host}
\subsection{Generic input}

\begin{frame}{Input data}
  Each row describes an observed interaction between two computers:
  \[
    e_i = (t_i, s_i, d_i, c_i).
  \]

  \begin{center}
    \small
    \begin{tabular}{@{}ll@{}}
      \toprule
      Required field & Meaning \\
      \midrule
      \code{timestamp} & $t_i$: event time in seconds \\
      \code{source_computer} & $s_i$: computer initiating the event \\
      \code{destination_computer} & $d_i$: destination computer \\
      \code{event} & $c_i$: categorical description of the event \\
      \bottomrule
    \end{tabular}
  \end{center}

  \small
  Optional fields include a unique \code{event_id}, \code{source_user},
  \code{success}, and an evaluation \code{label}.

  \note{The generic input is a table, usually a pandas DataFrame. The four
  required fields preserve both time and the direction of the interaction.
  Timestamps must be finite numeric seconds and host endpoints must be known.
  Keep a unique physical-row ID before partitioning the data. The security label
  is evaluation metadata. It does not define the event category or the graph.
  Source: deepcase_gnn/preprocessing/preprocessor.py, validate_frame.}
\end{frame}

\subsection{Historical host graphs}

\begin{frame}{Historical graph snapshots}
  For events at time $t$, collect earlier activity in a window of width $W$:
  \[
    \mathcal{H}_t = \{e_i : t-W \leq t_i < t\},
    \qquad G_t = (V_t, E_t).
  \]

  \begin{itemize}
    \item \textbf{Nodes:} computers in the historical window, plus the
      source and destination computers of current events.
    \item \textbf{Directed edges:} one edge $(s_i,d_i)$ for each distinct
      source--destination pair in $\mathcal{H}_t$.
    \item \textbf{Repeated interactions:} one edge per pair, with an
      interaction count and historical attributes.
  \end{itemize}

  \small
  All events at $t$ share the same snapshot. Their categories, edges, and
  statistics enter history only at later timestamps.

  \note{Build one graph per distinct timestamp at which events are scored.
  The window includes its lower bound and excludes time t, so simultaneous
  events cannot provide context to one another. The current endpoints can add
  isolated nodes, but current events do not create historical edges. Preserve
  direction: an edge from A to B does not automatically create B to A.
  The implementation constructs snapshots lazily. Its default window is 300
  seconds, which is a configuration choice rather than an evaluated optimum.
  Source: deepcase_gnn/preprocessing/preprocessor.py, HostGraphDataset.__getitem__.}
\end{frame}

\begin{frame}{Host and interaction features}
  \begin{columns}[T,onlytextwidth]
    \begin{column}{0.49\textwidth}
      \textbf{For each computer}
      \begin{itemize}
        \item Last $L$ outgoing event categories, in chronological order.
        \item Counts of outgoing events, distinct source users, successes,
          and failures.
      \end{itemize}
    \end{column}
    \begin{column}{0.47\textwidth}
      \textbf{For each directed pair}
      \begin{itemize}
        \item Interaction count and distinct source-user count.
        \item Failure fraction and age of the latest interaction.
      \end{itemize}
    \end{column}
  \end{columns}

  \medskip
  \small
  Fit the category vocabulary on training data only.
  Reserve token 0 for padding and token 1 for unseen categories.

  \medskip
  Counts and pair attributes use the whole window, even when a host's
  sequence exceeds $L$ events.

  \note{A host sequence contains only events for which that host is the source.
  Retain the last L categories, then right-pad shorter sequences with zeros.
  A destination-only host has an empty outgoing sequence and zero outgoing
  statistics. User and outcome statistics require the corresponding optional
  input columns. Node counts and pair counts use log(1+count). Pair recency is
  (t - latest pair time)/W. The exact formulas appear in the backup slide.
  The default sequence cap is 10 events. Hosts have no trainable identity
  embeddings. Standard GraphSAGE uses graph topology, while the prediction head
  also receives the current pair's historical attributes. Stored edge weights
  and edge attributes do not weight the GraphSAGE aggregation in this version.
  Sources: preprocessing/preprocessor.py and context_builder/context_builder.py
  inside deepcase_gnn.}
\end{frame}

\subsection{LANL example}

\begin{frame}{LANL authentication input}
  LANL authentication records contain nine fields. Computers and accounts
  use consistent anonymous identifiers, and time counts seconds from the
  beginning of collection.

  \medskip
  \textbf{Illustrative authentication record}
  \begin{columns}[T,onlytextwidth]
    \begin{column}{0.51\textwidth}
      \small
      \begin{tabular}{@{}ll@{}}
        \toprule
        Field & Value \\
        \midrule
        Time & \code{120} \\
        Source user & \code{U7@DOM1} \\
        Destination user & \code{U7@DOM1} \\
        Source computer & \code{C10} \\
        Destination computer & \code{C20} \\
        \bottomrule
      \end{tabular}
    \end{column}
    \begin{column}{0.47\textwidth}
      \small
      \begin{tabular}{@{}ll@{}}
        \toprule
        Field & Value \\
        \midrule
        Authentication & \code{Kerberos} \\
        Logon type & \code{Network} \\
        Orientation & \code{LogOn} \\
        Result & \code{Success} \\
        \bottomrule
      \end{tabular}
    \end{column}
  \end{columns}

  \medskip
  \footnotesize
  Source: \href{https://csr.lanl.gov/data/cyber1/}{LANL dataset documentation}.
  This implementation reads \code{auth.txt.gz} for graph construction.

  \note{LANL publishes several event sources. This implementation uses the
  authentication file for graph construction, with the optional Red Team file
  for evaluation. Do not imply that process, DNS, or flow files enter the current
  graph. The displayed values form a synthetic teaching example in the LANL
  schema, not an actual quoted authentication row. A computer identifier denotes
  an anonymized host, not a hardware type. Dataset source:
  https://csr.lanl.gov/data/cyber1/. Adapter source:
  deepcase_gnn/preprocessing/lanl.py, AUTH_COLUMNS and read_lanl.}
\end{frame}

\begin{frame}{LANL normalization}
  The adapter preserves the timestamp and both host endpoints. It builds
  an event category from:
  \[
    \begin{aligned}
      c_i = (&\text{authentication type},\ \text{logon type},\\
             &\text{orientation},\ \text{result},\ \text{user mismatch}).
    \end{aligned}
  \]

  \small
  In the worked example, abbreviate two categories as:
  \[
    \begin{aligned}
      S &= (\text{Kerberos},\text{Network},\text{LogOn},\text{Success},\text{False}),\\
      F &= (\text{NTLM},\text{Network},\text{LogOn},\text{Fail},\text{False}).
    \end{aligned}
  \]

  \begin{itemize}
    \item Default scope: \code{LogOn} with distinct, known source and
      destination computers.
    \item Keep all authentication mechanisms, account categories, and
      success/failure outcomes within that scope.
  \end{itemize}

  \note{User mismatch compares the full source and destination account strings.
  It is False for the illustrative U7@DOM1 to U7@DOM1 event. The adapter stores
  the five category fields as a compact JSON array of strings, including the
  literal string False or True. The preprocessor later maps that signature to
  a token using only the training vocabulary. S and F are slide shorthand, not
  the literal event strings or learned token numbers. Host and raw user IDs are
  not category vocabulary entries. The default remote-logon filter removes
  local events, other orientations, and unknown or empty endpoints. It does
  not merge duplicate input rows. An all-events configuration also exists.
  Source: deepcase_gnn/preprocessing/lanl.py, EVENT_FIELDS and read_lanl.}
\end{frame}

\begin{frame}{LANL window example}
  \textbf{Illustrative settings:} $t=200$ s and $W=100$ s.
  The graph uses the interval $[100,200)$.

  \smallskip
  \small All rows use \code{U7@DOM1} as both user accounts.

  \begin{center}
    \small
    \begin{tabular}{@{}rllll@{}}
      \toprule
      Time & Source & Destination & Category & Role at $t=200$ \\
      \midrule
       90 & \code{C30} & \code{C10} & $S$ & Outside window \\
      120 & \code{C10} & \code{C20} & $S$ & History \\
      160 & \code{C10} & \code{C20} & $F$ & History \\
      180 & \code{C20} & \code{C30} & $S$ & History \\
      \midrule
      200 & \code{C30} & \code{C10} & $F$ & Current target \\
      \bottomrule
    \end{tabular}
  \end{center}

  \begin{itemize}
    \item The three historical rows create two directed edges.
    \item The two \code{C10} to \code{C20} interactions share one edge.
    \item The current event does not create a \code{C30} to \code{C10}
      edge in this snapshot.
  \end{itemize}

  \note{Walk through the table by role. The event at time 90 has expired from
  the window. Times 120, 160, and 180 provide historical context. The target
  at time 200 selects its source and destination representations, but its
  category F is a prediction target. If other events also occurred at time
  200, they would share this graph without entering one another's context.
  Assume every row uses U7@DOM1 as both source and destination user, so mismatch
  is False and each active pair has one distinct source user. These are
  explanatory settings, not the implementation defaults. Source for window
  and aggregation rules: deepcase_gnn/preprocessing/preprocessor.py.}
\end{frame}

\begin{frame}{LANL graph and host histories}
  The historical graph at $t=200$ has three computer nodes:
  \[
    \mathrm{C10}
      \xrightarrow{\text{2 events}}
    \mathrm{C20}
      \xrightarrow{\text{1 event}}
    \mathrm{C30}.
  \]

  With $L=3$, the outgoing host sequences and counts are:
  \begin{center}
    \small
    \begin{tabular}{@{}l@{\qquad}l@{\qquad}l@{}}
      \toprule
      Host & Sequence & Counts $(n,u,s,f)$ \\
      \midrule
      \code{C10} & $[S,F,0]$ & $(2,1,1,1)$ \\
      \code{C20} & $[S,0,0]$ & $(1,1,1,0)$ \\
      \code{C30} & $[0,0,0]$ & $(0,0,0,0)$ \\
      \bottomrule
    \end{tabular}
  \end{center}

  \small
  $S$ and $F$ stand for category tokens, and $0$ denotes padding.
  Counts represent outgoing events, distinct source users, successes, and
  failures. Node features apply $\log(1+\text{count})$.

  \medskip
  The target pair $(\mathrm{C30},\mathrm{C10})$ has no edge in the window,
  so its historical pair attributes are zero.

  \note{The graph captures relationships across computers, while each sequence
  captures one computer's outgoing behavior. C30 has no outgoing history in
  this window, even though it is a destination and the source of the current
  target. Its historical sequence and outgoing statistics are therefore zero.
  Incoming GraphSAGE messages can still provide context from C20. The sequence
  entries S and F stand for their training-vocabulary tokens, not literal model
  inputs. The edge labels display stored interaction counts. Standard SAGEConv
  does not use those counts as aggregation weights in the current implementation.
  This graph and the encoded host histories become inputs to the GRU and GNN.
  Sources: deepcase_gnn/preprocessing/preprocessor.py and
  deepcase_gnn/context_builder/context_builder.py.}
\end{frame}

\appendix

\begin{frame}{Backup: feature definitions and evaluation labels}
  \small
  For a host $v$, let $n_v,u_v,s_v,f_v$ count its outgoing events,
  distinct source users, successes, and failures:
  \[
    x_v = [\log(1+n_v),\log(1+u_v),\log(1+s_v),\log(1+f_v)].
  \]

  For an ordered pair $(a,b)$, let $n_{ab}$ count events,
  $u_{ab}$ count distinct source users, and $\rho_{ab}$ be its failure fraction:
  \[
    a_{ab} = \left[\log(1+n_{ab}),\log(1+u_{ab}),\rho_{ab},
      \frac{t-t^{\mathrm{last}}_{ab}}{W}\right].
  \]

  \begin{itemize}
    \item GraphSAGE uses topology. The event head uses the target pair's
      historical attributes.
    \item Optional Red Team labels match
      $(\text{time},\text{source user},\text{source host},\text{destination host})$
      exactly and serve evaluation only.
    \item Label 0 means no known Red Team match.
  \end{itemize}

  \note{All feature counts use the full historical window. Pair failure
  fraction is the number of Fail rows divided by all rows for that pair.
  Recency is normalized by window width W. Stored edge_weight is the raw
  interaction count, but neither edge_weight nor edge_attr is passed into
  SAGEConv. The event head receives the historical attributes for its selected
  pair, or zeros for an unseen pair. A security label is separate from the event
  category being predicted and does not enter model training or features.
  A missing Red Team match does not prove that an event was benign.
  Sources: deepcase_gnn/preprocessing/preprocessor.py,
  deepcase_gnn/preprocessing/lanl.py, and
  deepcase_gnn/context_builder/context_builder.py. Dataset label source:
  https://csr.lanl.gov/data/cyber1/.}
\end{frame}

\end{document}
